% ==================================================
% The Binomial Theorem
% Discrete Calculus — Foundations
% ==================================================

\section{The Binomial Theorem}
\label{sec:binomial-theorem}

\begin{tcolorbox}[colback=thmbox, colframe=thmborder, arc=2pt,
  left=6pt, right=6pt, top=4pt, bottom=4pt,
  title={\small\textbf{Theorem (Binomial Theorem)}},
  fonttitle=\small\bfseries]
For any integer \(n \ge 0\) and real numbers \(x,y\),
\[
(x+y)^n
=
\sum_{k=0}^{n} \binom{n}{k} x^{\,n-k} y^{\,k},
\]
where the \emph{binomial coefficients} are
\[
\binom{n}{k}
=
\frac{n!}{k!(n-k)!}.
\]
\end{tcolorbox}
\label{thm:binomial}

\begin{remark*}[Fully quantified form]
\(\forall n \in \mathbb{N},\; \forall x,y \in \mathbb{R}: \displaystyle(x+y)^n = \sum_{k=0}^{n}\binom{n}{k}x^{n-k}y^{k}\).
\end{remark*}

\begin{remark*}[Combinatorial meaning]
The binomial coefficients count the number of ways to choose \(k\) factors of
\(y\) from the \(n\) factors in the product \((x+y)(x+y)\cdots(x+y)\).
Thus \(\binom{n}{k}\) is the number of ways to select \(k\) objects from a
set of \(n\), and the theorem describes how a power of a sum expands into
a sum of products.
\end{remark*}

\section*{Basic Properties of Binomial Coefficients}

\begin{align*}
\binom{n}{0} &= 1, &
\binom{n}{n} &= 1, &
\binom{n}{k} &= \binom{n}{n-k}.
\end{align*}

They also satisfy the \emph{Pascal recursion}:
\[
\binom{n+1}{k}
=
\binom{n}{k}+\binom{n}{k-1}.
\]
This recursion produces Pascal's triangle, in which each entry is the
sum of the two entries directly above it.

\begin{center}
\begin{tikzpicture}[scale=1]
\foreach \n in {0,...,5}
{
  \foreach \k in {0,...,\n}
  {
    \node at (\k-\n/2,-\n) {$\binom{\n}{\k}$};
  }
}
\end{tikzpicture}
\end{center}

\begin{remark*}[Pascal's triangle]
Each entry \(\binom{n}{k}\) is obtained from the sum of the two entries
above it: \(\binom{n}{k}=\binom{n-1}{k-1}+\binom{n-1}{k}\).
These coefficients appear throughout discrete calculus, especially in
formulas for higher finite differences and in Newton's forward difference
expansion.
\end{remark*}

\section*{Connection with Discrete Calculus}

The binomial theorem plays a central role in discrete calculus because
difference operators naturally produce binomial coefficients.

\begin{proposition}[Higher difference formula]
\label{prop:higher-diff-formula}
For any function \(f\),
\[
\Delta^n f(x)
=
\sum_{k=0}^{n} (-1)^{\,n-k}
\binom{n}{k}
f(x+k).
\]
\end{proposition}

\begin{remark*}[Fully quantified form]
\(\forall n \in \mathbb{N},\; \forall f : \mathbb{Z}\to\mathbb{R},\; \forall x \in \mathbb{Z}:\;
\Delta^n f(x) = \displaystyle\sum_{k=0}^{n}(-1)^{n-k}\binom{n}{k}f(x+k)\).
\end{remark*}

\begin{remark*}[Proof idea]
Applying the difference operator \(\Delta = E - I\) (where \(E\) is the shift
operator) repeatedly gives \(\Delta^n = (E-I)^n\). Expanding by the binomial
theorem yields the coefficients \((-1)^{n-k}\binom{n}{k}\), and \(E^k f(x) = f(x+k)\).
This derivation is carried out in full in the Difference Operators section.
\end{remark*}

\section*{Explicit Examples}

Applying the difference operator repeatedly to a function \(f\) yields
\[
\begin{aligned}
\Delta f(x) &= f(x+1)-f(x), \\
\Delta^2 f(x) &= f(x+2)-2f(x+1)+f(x), \\
\Delta^3 f(x) &= f(x+3)-3f(x+2)+3f(x+1)-f(x).
\end{aligned}
\]
The coefficients \(1, 2, 1\) and \(1, 3, 3, 1\) are precisely the rows of
Pascal's triangle, confirming that binomial coefficients determine the pattern
of coefficients in higher-order finite differences.
